\documentclass[11pt,a4paper]{article}
\usepackage{fontspec,amsmath,amssymb,amsthm,microtype}
\setmainfont{lmroman10-regular.otf}[BoldFont=lmroman10-bold.otf,ItalicFont=lmroman10-italic.otf,BoldItalicFont=lmroman10-bolditalic.otf]
\setmonofont{lmmono10-regular.otf}[Scale=MatchLowercase]
\usepackage[margin=25mm]{geometry}
\usepackage[unicode,hidelinks]{hyperref}
\setlength{\emergencystretch}{3em}
\setlength{\parskip}{0.4em}
\setlength{\parindent}{0pt}
\allowdisplaybreaks[2]
\raggedbottom
\title{\texorpdfstring{Interacting Particle Langevin under LSI}{Interacting Particle Langevin under LSI}}
\author{LangevinLean mathematical exposition}
\date{}
\begin{document}\maketitle
The accompanying HTML pairs this mathematical exposition with the exact Lean statements and records the scope of the formalization.
The lifted-target entropy bound follows from the full actual Gaussian ULA proof and finite initial joint KL. The LSI hypothesis concerns the actual lifted target; all interpolation and moment estimates are derived.
\section{\texorpdfstring{Theorem}{Theorem}}


Let \(N\ge1\), \(D=d_\theta+Nd_x\), assume \(D\ge1\), and let the joint potential \(U(\theta,x)\) have an \(L\)-Lipschitz gradient with \(L>0\). In scaled coordinates set


\[z=(\theta,x^1/\sqrt N,\ldots,x^N/\sqrt N),\qquad
V_N(z)=\sum_{i=1}^NU(\theta,\sqrt N z_i),\quad
\Pi_N(dz)\propto e^{-V_N(z)}dz.\]


Assume an LSI with constant \(\beta_N>0\) for this actual lifted target. Compatibility gives \(\beta_N\le NL\) by the positive-dimensional argument on the shared ULA page; the final Lean theorem derives this compatibility internally. For finite initial joint entropy and


\[0<\gamma\le\frac{\beta_N}{4NL^2},\]


the IPLA laws \(\mu_k\), equivalently ULA for \(V_N\) at step \(\gamma/N\), satisfy


\[\frac{H_{\Pi_N}(\mu_k)}N\le
e^{-\beta_N\gamma k/N}\frac{H_{\Pi_N}(\mu_0)}N
+\frac{8\gamma DL^2}{\beta_N}.\]


In particular, if the lifted target has the admissible constant \(\beta_N=N\alpha\), then


\[\boxed{\frac{H_{\Pi_N}(\mu_k)}N\le
e^{-\alpha\gamma k}\frac{H_{\Pi_N}(\mu_0)}N
+\frac{8\gamma L^2}{\alpha}\left(d_x+\frac{d_\theta}N\right).}\]


The condition is now \(0<\gamma\le\alpha/(4L^2)\). An LSI for the original pair target is not substituted for the lifted-target hypothesis. The Lean result derives all interpolation estimates from the shared ULA theorem and finite initial joint KL. If \(D=0\), both laws are the unique point mass and their entropy is zero.


\section{\texorpdfstring{The geometry fixes the dimension, gradient constant, and step}{The geometry fixes the dimension, gradient constant, and step}}


Each physical pair \((\theta,x^i)\) is extracted from the scaled state by \((\theta,\sqrt N z_i)\). Differentiating the lifted sum gives the averaged parameter drift and ordinary latent drifts at step \(\gamma/N\). The corresponding Euclidean state has


\[D=d_\theta+Nd_x,\qquad L_N=NL,\qquad h_N=\gamma/N.\]


The parameter block of the lifted gradient is a sum over particles; each latent block contributes a factor \(\sqrt N\). Squaring and summing the blocks, then applying Cauchy–Schwarz to the common-parameter sum, proves the gradient Lipschitz constant \(NL\). The coordinate Gaussian laws and independence identify the actual noise-driven IPLA law with this ULA transition. No equality between \(d_\theta\) and \(d_x\) is needed.



\textbf{Formal scope.} The lifted geometry and identification with the actual Gaussian particle algorithm are checked.



{\raggedright\noindent Lean declarations: \texttt{particle\allowbreak{}Laws\_\allowbreak{}eq\_\allowbreak{}laws}.\par}

\section{\texorpdfstring{Apply the ULA one-\allowbreak{}step estimate with the lifted target}{Apply the ULA one-step estimate with the lifted target}}


Apply the ULA entropy argument to \((V_N,\Pi_N)\), with constant \(\beta_N\), dimension \(D\), smoothness \(NL\), and step \(\gamma/N\). Its step restriction becomes


\[\frac\gamma N\le\frac{\beta_N}{4(NL)^2}
\quad\Longleftrightarrow\quad
\gamma\le\frac{\beta_N}{4NL^2}.\]


The one-step forcing simplifies without losing a particle factor:


\[6h_N^2D L_N^2=6(\gamma/N)^2D(NL)^2=6\gamma^2DL^2.\]


Consequently the unnormalized joint entropy obeys


\[H_{k+1}\le e^{-\beta_N\gamma/N}H_k+6\gamma^2DL^2.\]


The Lean proof constructs probability and finite-KL guards first, applies the shared one-step theorem, and then rewrites the rescaled constants. Every analytic input is supplied by the actual Gaussian ULA construction.



\textbf{Formal scope.} The particle substitution and scalar algebra use the shared ULA proof from smoothness, LSI, and finite initial KL.



{\raggedright\noindent Lean declarations: \texttt{ula\_\allowbreak{}kl\_\allowbreak{}one\_\allowbreak{}step}.\par}

\section{\texorpdfstring{Normalize by the particle number only after deriving the joint bound}{Normalize by the particle number only after deriving the joint bound}}


Iteration gives the unnormalized floor


\[\frac{8h_NDL_N^2}{\beta_N}
=\frac{8\gamma NDL^2}{\beta_N}.\]


Divide the entire inequality by positive \(N\), including the initial entropy and the floor. This proves


\[\frac{H_k}N\le e^{-\beta_N\gamma k/N}\frac{H_0}N
+\frac{8\gamma DL^2}{\beta_N}.\]


Now, and only under the separate lifted-target assumption \(\beta_N=N\alpha\), use \(D/N=d_x+d_\theta/N\). The rate is \(e^{-\alpha\gamma k}\), and the per-particle floor is \((8\gamma L^2/\alpha)(d_x+d_\theta/N)\). Shared parameter coordinates prevent a direct tensorization argument from a pair LSI; the target for which LSI is assumed remains \(\Pi_N\).



\textbf{Formal scope.} The general lifted-constant bound, division by N, and linear-LSI specialization are checked with their assumptions visible.



{\raggedright\noindent Lean declarations: \texttt{particle\_\allowbreak{}kl\_\allowbreak{}bound\_\allowbreak{}normalized\_\allowbreak{}of\_\allowbreak{}finite\_\allowbreak{}kl}, \texttt{particle\_\allowbreak{}kl\_\allowbreak{}bound\_\allowbreak{}linear\_\allowbreak{}lsi\_\allowbreak{}of\_\allowbreak{}finite\_\allowbreak{}kl}.\par}

\section{\texorpdfstring{Project to the parameter marginal using data processing}{Project to the parameter marginal using data processing}}


Let \(P_\theta\) be the parameter projection, \(\mu_k^\theta=P_{\theta\#}\mu_k\), and \(\Pi_N^\theta=P_{\theta\#}\Pi_N\). Relative entropy decreases under a measurable map:


\[H_{\Pi_N^\theta}(\mu_k^\theta)\le H_{\Pi_N}(\mu_k).\]


This controls the unnormalized parameter KL. In particular, under \(\beta_N=N\alpha\),


\[H_{\Pi_N^\theta}(\mu_k^\theta)
\le e^{-\alpha\gamma k}H_{\Pi_N}(\mu_0)
+\frac{8\gamma NL^2}{\alpha}\left(d_x+\frac{d_\theta}N\right).\]


The initial entropy remains the full joint entropy. Data processing does not allow it to be replaced by the smaller initial parameter entropy on the right. Likewise, a normalized joint bound must be multiplied by \(N\) before it bounds the unscaled parameter KL. The Lean proof establishes finite joint KL, transfers finiteness to the marginal, and only then compares the real values.



\textbf{Formal scope.} The parameter data-processing argument and all normalization factors are checked. This conclusion concerns a probability marginal, not distance to a likelihood maximizer.



{\raggedright\noindent Lean declarations: \texttt{parameter\_\allowbreak{}kl\_\allowbreak{}bound\_\allowbreak{}of\_\allowbreak{}finite\_\allowbreak{}kl}.\par}

\end{document}
